\section{Cramer's rule}

Cramer's rule is a method for solving certain square systems of linear equations using determinants. It is not the most efficient method for large systems, but it is conceptually important because it shows a direct connection between determinants and solvability.

The method applies to systems where the number of equations equals the number of unknowns and the coefficient matrix has non-zero determinant.

\subsection*{The two by two case}

Consider the system
\[
\begin{cases}
a x+b y=e,\\
c x+d y=f.
\end{cases}
\]
The coefficient matrix is
\[
A=\begin{pmatrix}
a&b\\
c&d
\end{pmatrix},
\]
and its determinant is
\[
\det A=ad-bc.
\]
If \(\det A\ne0\), the system has exactly one solution.

Cramer's rule says that the solution can be computed by replacing one column of the coefficient matrix at a time by the right-hand side column.

Define
\[
A_x=\begin{pmatrix}
e&b\\
f&d
\end{pmatrix},
\qquad
A_y=\begin{pmatrix}
a&e\\
c&f
\end{pmatrix}.
\]
Then
\[
x=\frac{\det A_x}{\det A},
\qquad
 y=\frac{\det A_y}{\det A}.
\]

\begin{remarkbox}
The denominator \(\det A\) is not a technical detail. If \(\det A=0\), the coefficient matrix is singular, and Cramer's rule does not apply.
\end{remarkbox}

\subsection*{A worked example}

\begin{examplebox}
Solve the system
\[
\begin{cases}
2x+y=5,\\
x-y=1.
\end{cases}
\]
The coefficient matrix is
\[
A=\begin{pmatrix}
2&1\\
1&-1
\end{pmatrix}.
\]
Its determinant is
\[
\det A=2\cdot(-1)-1\cdot1=-2-1=-3.
\]
Since \(\det A\ne0\), Cramer's rule applies.

For \(x\), replace the first column by the right-hand side:
\[
A_x=\begin{pmatrix}
5&1\\
1&-1
\end{pmatrix}.
\]
Then
\[
\det A_x=5\cdot(-1)-1\cdot1=-5-1=-6.
\]
So
\[
x=\frac{\det A_x}{\det A}=\frac{-6}{-3}=2.
\]

For \(y\), replace the second column:
\[
A_y=\begin{pmatrix}
2&5\\
1&1
\end{pmatrix}.
\]
Then
\[
\det A_y=2\cdot1-5\cdot1=2-5=-3.
\]
So
\[
y=\frac{\det A_y}{\det A}=\frac{-3}{-3}=1.
\]
Thus
\[
(x,y)=(2,1).
\]
\end{examplebox}

This result should be checked by substitution. We have
\[
2\cdot2+1=5,
\qquad
2-1=1.
\]
Both equations are satisfied.

\subsection*{The general square case}

For a square system
\[
Ax=b,
\]
where \(A\) is an \(n\times n\) matrix and \(b\) is a column vector, Cramer's rule applies if
\[
\det A\ne0.
\]
For each unknown \(x_i\), form the matrix \(A_i\) by replacing the \(i\)-th column of \(A\) with \(b\). Then
\[
x_i=\frac{\det A_i}{\det A}.
\]

The rule is elegant, but the number of determinant computations grows quickly. For this reason, Gaussian elimination is usually more practical for larger systems.

\subsection*{Why the determinant condition matters}

The condition
\[
\det A\ne0
\]
means that the coefficient matrix is invertible. Geometrically and structurally, the equations contain enough independent information to determine one solution.

If
\[
\det A=0,
\]
then the matrix is singular. This does not automatically mean that the system has no solution. It means that the system cannot have a unique solution determined by Cramer's rule. The system must be analyzed further: it may have no solution or infinitely many solutions.

\begin{examplebox}
Consider
\[
\begin{cases}
x+y=2,\\
2x+2y=4.
\end{cases}
\]
The coefficient matrix is
\[
A=\begin{pmatrix}
1&1\\
2&2
\end{pmatrix},
\]
and
\[
\det A=1\cdot2-1\cdot2=0.
\]
Cramer's rule does not apply. But the system is not inconsistent. The second equation is twice the first, so there are infinitely many solutions.
\end{examplebox}

This example is important. A zero determinant means loss of uniqueness, not automatically loss of existence.

\begin{summarybox}
When using Cramer's rule, ask:
\begin{itemize}
\item Is the system square?
\item What is the coefficient matrix \(A\)?
\item Is \(\det A\ne0\)?
\item Which column must be replaced to compute each unknown?
\item Does the final answer satisfy the original equations?
\item If \(\det A=0\), what further classification is needed?
\end{itemize}
\end{summarybox}

Cramer's rule shows how determinants can solve systems, but it must be used only when its structural condition is satisfied.